\chapter{Ausblick}
\label{chap:ausblick}

Die Ergebnisse dieser Arbeit legen mehrere Richtungen für weiterführende Forschung nahe, insbesondere im Hinblick auf die Übertragbarkeit modular trainierter Agenten in komplexeren Umgebungen, die Verbesserung der zugrundeliegenden Modellarchitektur sowie die Automatisierung des Lernprozesses.

\section{Multi-Agenten}

Ein zentrales Motiv dieser Arbeit bestand darin, Wege zu finden, wie stabile Lernprozesse in komplexen Multi-Agenten-Umgebungen ermöglicht werden können. Die initialen Experimente zeigten, dass ein direktes gemeinsames Training mehrerer Agenten zu instabilen und schwer interpretierbaren Verhaltensmustern führte. Durch die Reduktion auf isolierte Single-Agenten-Aufgaben konnten zentrale Teilfähigkeiten wie Navigation, Exploration und Flucht gezielt trainiert und analysiert werden. Die im Rahmen dieser modularen Trainingsstrategie gewonnenen Erkenntnisse bilden nun eine fundierte Ausgangsbasis, um das ursprüngliche Ziel unter kontrollierten Bedingungen erneut zu verfolgen.

Auf dieser Grundlage erscheint ein erneuter Versuch im Multi-Agenten-Kontext sinnvoll, bei dem nicht alle Agenten gleichzeitig trainiert werden, sondern vorab spezialisierte Komponenten einzeln geschult und anschließend gezielt in eine gemeinsame Umgebung integriert werden. Dies würde es erlauben, stabil trainierte Verhaltensmuster auf koordinierte oder konkurrierende Szenarien zu übertragen, ohne die bekannten Probleme ungeplanter Policy-Interaktion von Beginn an in Kauf zu nehmen.

Ein Schwerpunkt zukünftiger Untersuchungen könnte auf der Frage liegen, inwieweit sich die im Single-Agent-Setup erzielte Robustheit auch im Zusammenspiel mehrerer Agenten aufrechterhalten lässt. Dabei wäre insbesondere zu prüfen, ob ein schrittweises Curriculum Learning im Multi-Agenten-Training hilfreich sein kann, um konflikthafte Verhaltensmuster zu vermeiden und ein abgestimmtes Zusammenspiel zu fördern. Ergänzend könnten Methoden wie adaptive Policy-Auswahl oder modulare Architekturen zum Einsatz kommen, um situativ zwischen erlernten Verhaltenskomponenten zu wechseln und so flexibler auf komplexe Interaktionen zu reagieren.

\section{Modellarchitekturen: Attention und Embedding}

In dieser Arbeit wurde eine rekurrente Policy-Struktur auf Basis von LSTM-Modulen verwendet. Zukünftig wäre zu untersuchen, ob neuere Architekturen wie Attention-basierte Modelle Vorteile in Bezug auf Generalisierung und Kontextsensitivität bieten. Beispielsweise könnten \textit{Transformers} oder verwandte Mechanismen eingesetzt werden, um die Ray-basierte Beobachtungsstruktur effizienter auszuwerten, da sie in der Lage sind, langfristige Abhängigkeiten zu modellieren und relevante Umgebungsmerkmale selektiv zu gewichten \cite{parisotto2020stabilizing}.

Darüber hinaus eröffnet der Einsatz kontinuierlich trainierter Objekt-Embeddings neue Möglichkeiten im Bereich des \textit{Representation Learning}. Anstelle diskreter Klassifikationen könnten visuell ähnliche oder funktional verwandte Objekte durch gemeinsame Repräsentationen beschrieben werden. Dies würde insbesondere dann Vorteile bringen, wenn sich die Umgebung verändert, dabei jedoch grundlegende strukturelle Eigenschaften erhalten bleiben. Die Arbeit von Baker et al. zeigt exemplarisch, wie Objekt-Embeddings in komplexen Multi-Agenten-Umgebungen erfolgreich eingesetzt werden können, um funktionale Zusammenhänge zwischen Objekten zu erfassen und generalisierbares Verhalten zu fördern \cite{baker2020emergenttoolusemultiagent}. Zu berücksichtigen ist allerdings, dass mit wachsender Modellkapazität auch neue Herausforderungen entstehen. Insbesondere besteht das Risiko einer Überanpassung an lokale Trainingsmuster. Daher bleibt eine ausgewogene Abstimmung zwischen Modellkomplexität und Verallgemeinerungsfähigkeit ein zentraler Aspekt zukünftiger Architekturentscheidungen.

\section{Automatisierung des Lernprozesses}

Neben der Modellerweiterung bietet auch die Strukturierung des Lernprozesses selbst Potenzial für Verbesserungen. In dieser Arbeit wurden Subtasks manuell definiert und in einer festen Reihenfolge trainiert. Eine interessante Erweiterung wäre der Einsatz von \textit{automatischem Curriculum Learning}, bei dem die Schwierigkeit der Trainingsumgebung adaptiv an den Lernfortschritt des Agenten angepasst wird. Dies würde dem Agenten erlauben, komplexe Fähigkeiten schrittweise aufzubauen, ohne dass explizit definierte Trainingsphasen nötig sind. Darüber hinaus könnten alternative Reinforcement-Learning-Methoden – etwa Off-Policy-Algorithmen wie SAC oder DDPG – für bestimmte Subtasks vorteilhaft sein. Auch hybrid modellbasierte Ansätze wären denkbar, insbesondere um Lerndaten effizienter zu nutzen oder Planungskomponenten zu integrieren.
